\chapter{Obesity Screening}
\section*{What Is This Test?} Obesity screening estimates excess adiposity and related health risk using body measurements and clinical assessment.\section*{Alternative Names} BMI screening; weight-status assessment; obesity evaluation.\section*{Principle} BMI, waist circumference, growth charts, and health history estimate risk but do not measure body fat perfectly.\section*{Why This Test Is Ordered} It identifies cardiometabolic risk and guides prevention, nutrition, activity, and treatment discussions.\section*{Primary vs Secondary Test} Measurement and history are primary; glucose, lipids, liver tests, sleep studies, or endocrine tests are targeted secondary studies.\section*{How This Test Is Performed} Height, weight, waist, blood pressure, medicines, conditions, diet, activity, and sleep are reviewed.\section*{What Preparation Is Needed} Wear light clothing and remove shoes when measured; no fasting is needed for the assessment itself.\section*{Understanding Results} BMI is a screening measure and may misclassify muscular people, older adults, pregnancy, or different body compositions.\section*{Follow-up Tests} A1C or glucose, lipids, liver tests, sleep evaluation, and nutrition or behavioral assessment may follow.\section*{References} \begin{enumerate}\item MedlinePlus. Obesity Screening.\item U.S. Preventive Services Task Force. Screening for obesity in adults and children.\end{enumerate}\clearpage\section*{Understanding Results and Follow-up}
The laboratory report should identify the specimen, method, units, reference interval or decision limit, and whether the result is positive, negative, abnormal, or inconclusive. Reference intervals vary with age, sex, pregnancy, specimen type, assay, and laboratory; a result outside a range is not automatically a diagnosis, and a result within a range does not always exclude disease. Discuss unexpected results with the ordering clinician, who can decide whether repeat or confirmatory testing, treatment, or referral is appropriate. Do not change medicines or delay urgent care based on an isolated result.
\section*{Regulatory Status and Authorization Date}This chapter describes a test category, not a specific commercial product. FDA, CDSCO, EU IVDR, and MHRA approval or registration dates therefore cannot be assigned without the assay manufacturer, product name, instrument, and intended use. For laboratory-developed tests, an individual agency approval date may not exist. See the Preface for the interpretation of regulatory status.
\clearpage
